\subsection*{多层并行}

可以嵌套声明不同的并行层次：

\begin{lstlisting}[language=croqtile,caption={嵌套的两级并行}]
parallel p by 2 {
  // parallel-level-0
  parallel q by 12 {
    // parallel-level-1
  }
}
\end{lstlisting}

本例的两层范围分别为 \co@2@ 和 \co@12@，共有 24 种索引组合。\co@p@ 属于外层，\co@q@ 属于内层，每个组合标识一份工作。

层次划分也决定代码在哪个范围内执行，以及哪些工作需要协调。组织代码时，应让协作关系与并行层次相匹配。

两层之间没有其他语句时，也可以合写为一条声明：

\begin{lstlisting}[language=croqtile,caption={合写的两级并行}]
parallel p by 2, q by 12 { ... }
\end{lstlisting}

逗号分隔的写法与前面的嵌套声明表达相同的两级并行。
